\scanpage{28}
\margin{-4}
\uline{Exercice 3} : Calculer les intégrales indéfinies suivantes :
\[
\text{a) }\int x^5 dx\ ;\qquad
\text{b) }\int\frac1{x^3}\,dx;\qquad
\text{c) }\int x\sqrt{x}\,dx\ ;\qquad
\text{d) }\int\frac1{x\sqrt[3]{x}}\,dx
\]

\uline{Propriétés} :

Primitive du \uline{produit} d'une fonction par une constante : $\lambda f(x)$.

On fait le produit de l'intégrale par la constante : $\lambda\int f(x)\,dx$.

Primitive de la \uline{somme} de deux fonctions : $f(x)+g(x)$

ou de la \uline{différence} de deux fonctions : $f(x)-g(x)$

on fait la somme (ou la différence) des intégrales.

Soit :
\[
\begin{aligned}
\int\lambda f(x)\,dx&=\lambda\int f(x)\,dx\\
\int[f(x)+g(x)]\,dx&=\int f(x)\,dx+\int g(x)\,dx\\
\int[f(x)-g(x)]\,dx&=\int f(x)\,dx-\int g(x)\,dx
\end{aligned}
\]

\uline{Attention} ! on n'a pas de formule analogue pour le produit.

\uline{En résumé} : Etant donné deux fonctions $f$ et $g$ dont on connaît une primitive, respectivement $F$ et $G$, alors $\lambda F+\mu G$ est une primitive de la \uline{combinaison linéaire} $\lambda f+\mu g$.

\uline{Exercice 4} : Calculer les intégrales indéfinies suivantes :
\[
\text{a) }\int(3x^5-5x^2+1)\,dx\ ;\qquad
\text{b) }\int2(\sin x-\cos x)\,dx
\]

\uline{Transformations} :

Il est parfois nécessaire de transformer l'expression de la fonction dont on cherche une primitive,

- pour la simplifier,

- pour l'écrire comme une combinaison linéaire de fonctions dont on connaît une primitive.

\uline{Exemple} :
\[
\begin{aligned}
\int\frac{\sin^2x}{1+\cos x}\,dx
&=\int\frac{1-\cos^2x}{1+\cos x}\,dx
=\int\frac{(1-\cos x)(1+\cos x)}{1+\cos x}\,dx=\\
&=\int(1-\cos x)\,dx
=\int dx-\int\cos x\,dx=x-\sin x+c.
\end{aligned}
\]

\uline{Exercice 5} : En transformant l'expression donnée, calculer les intégrales indéfinies suivantes :
\[
\begin{gathered}
\text{a) }\int(x^2+1)^2dx\ ;\quad
\text{b) }\int\operatorname{tg}^2x\,dx\ ;\quad
\text{c) }\int\sin^2\frac x2\,dx\ ;\quad
\text{d) }\int\frac{2+x^2}{1+x^2}\,dx\\
\text{e) }\int(\sin^3x+\cos^2x\sin x)\,dx\ ;\quad
\text{f) }\int\sin\frac x2\cos\frac x2\,dx\ ;\quad
\text{g) }\int\frac{(x-1)(x+1)}x\,dx
\end{gathered}
\]

\uline{Faire l'autotest 1}
\margin{27}
